\chapter{Optimal Packing of Non-Convex Polygons}

This project is an excercise in optimization and High performance computing.
   It is inspired in the classical problem of sphere packing with a twist:
   We seek to find the smallest square container that can fit N identical
    non-convex polygons without overlap.

\section{Mathematical Formulation}

We consider a set of $N$ identical non-convex polygons
\[
\poly = \{P_1, \dots, P_N\},
\]
Each polygon can be described with the state variables $(x_i, y_i, \theta_i) \in \R^2 \times [0,2\pi)$.

\begin{definition}[Feasible Configuration]
A configuration $\conf$ is feasible if:
\begin{enumerate}
  \item All vertices lie within the square domain $[0,L]^2$.
  \item No two polygons overlap.
\end{enumerate}
\end{definition}

To objective is to find the minimal $L$ such that all polygons 
are placed inside the square without overlap. Formally, we seek to solve:
\[
\begin{aligned}
\min_{L, \conf} \quad & L \\
\text{s.t.} \quad & \conf \text{ is a feasible configuration of } \poly \text{ in } [0,L]^2.
\end{aligned}
\]


\subsection{Feasibility-Based Search}

A common strategy is to reduce the optimization problem to a sequence of feasibility tests. For a fixed container size $L$, one asks whether there exists a configuration $\conf$ such that all polygons fit inside $[0,L]^2$ without overlap. This leads to the decision problem
\[
\text{Find } \conf \quad \text{such that } \conf \text{ is feasible in } [0,L]^2.
\]

By performing a search over values of $L$. Using bisection,
 one can approximate the minimal feasible container size.
  This approach shifts the difficulty from direct optimization
   to repeated feasibility detection in a highly constrained configuration space.

\subsection{Energy-Based Formulations}

An alternative and widely used approach replaces hard geometric
 constraints with soft penalties by introducing an energy function of the form
\[
E(\conf;L) = E_{\text{overlap}}(\conf) + E_{\text{boundary}}(\conf;L),
\]
where overlaps and boundary violations contribute positive energy.
 In this formulation, feasible configurations correspond to zero-energy states.

This transforms the constrained optimization problem into an
 unconstrained (but non-convex) minimization problem over
  $\mathbb{R}^{2N} \times S^1$.
  The resulting energy landscape is typically rugged,
   containing many local minima corresponding to jammed or partially 
   overlapping configurations.

Energy-based formulations naturally motivate the use of stochastic optimization techniques,
 such as simulated annealing and Monte Carlo methods. We use them to explore the configuration space
  and find feasible or near-feasible configurations.
   which are designed to explore complex non-convex landscapes.


\paragraph{Computational and algorithmic challenges.}
The packing problem studied in this work presents several fundamental difficulties,
 both geometric and algorithmic.

\begin{enumerate}
\item \textbf{Highly non-convex energy landscape.}
The objective function associated with packing
design by associating a penalties to overlaps and boundary violations
is highly non-convex. Additionally, it exhibits a large number of local minima. 
As will be shown in the next chapter, even simplified model like 
the Ising model  possesses a 
 non-convex energy landscape. As a consequence,
  purely local optimization methods get trapped in suboptimal
   configurations, requiring use of global or 
   stochastic optimization techniques.

\item \textbf{Geometric complexity due to non-convex shapes.}
By the nature of non-convex polygons,
 the overlap checks require careful geometric
  computation, often involving decomposition into convex components
   or the evaluation of separating axes. This geometric complexity 
   significantly increases the computational cost of evaluating
    candidate configurations.

\item \textbf{Exponential growth of the search space.}
The configuration space grows exponentially with the number of objects $N$,
 as each polygon contributes translational and rotational degrees of freedom.
  Consequently, exhaustive search methods become infeasible even for moderately
   large systems.

\item \textbf{Computational intractability (NP-hardness).}
The problem is computationally intractable in general.
 This follows from the simplified version being hard: 
 if all polygons are restricted to be circles. 
 This can be seen by considering the associated \emph{decision problem}: 
 given a set of circles with prescribed radii and a square container of 
 fixed side length, determine whether all circles can be placed inside the 
 square without overlap.

This decision problem is known as the \emph{circle packing problem}
 and has been shown to be NP-complete. Since circle packing is a special case of the more general polygon packing problem, it follows that the general problem is NP-hard.

\end{enumerate}

